\label{sec:presa_ramjet}
\section{Presa d'Aria Supersonica}

% VERSIONE ACCORCIATA
La presa d'aria del ramjet è stata dimensionata nel punto di progetto
($M_0=3.5$, $h=25000$ m) con l'obiettivo di decelerare il flusso fino a
condizioni compatibili con la combustione, massimizzando il recupero di
pressione totale e contenendo al tempo stesso l'ingombro del componente.

\paragraph{Architettura: presa a compressione esterna multirampa.}
È stata adottata una configurazione a compressione esterna su tre rampe,
ciascuna generatrice di un urto obliquo, seguite da un urto normale
terminale che porta il flusso a condizioni subsoniche. Tale architettura
è stata scelta per limitare le elevate perdite di pressione totale che,
a $M_0=3.5$, caratterizzerebbero una decelerazione mediante un singolo
urto normale. La configurazione a tre rampe rappresenta quindi un
compromesso tra recupero di pressione totale, complessità geometrica e
ingombro della presa. Il flusso è stato modellato come adiabatico e non
viscoso, assumendo $T_0=\mathrm{cost}$ e considerando esclusivamente le
perdite di pressione totale associate agli urti, trascurando gli effetti
viscosi e di strato limite sulle rampe.

\paragraph{Ottimizzazione degli urti.}
Definito il numero di rampe, la deflessione complessiva è stata ripartita
secondo il criterio di Oswatitsch. A parità di condizioni iniziali, il
recupero di pressione totale viene massimizzato distribuendo la
compressione tra urti di uguale intensità. Nel modello è stato quindi
imposto lo stesso rapporto di pressione statica attraverso i tre urti
obliqui e tale rapporto è stato determinato mediante la massimizzazione
del recupero di pressione totale complessivo. Ne risultano deflessioni
progressivamente crescenti, pari a $11.8^\circ$, $14.4^\circ$ e
$17.7^\circ$.

A valle del terzo urto il flusso permane supersonico; è stato quindi
introdotto un urto normale terminale per raggiungere condizioni
subsoniche. Questa soluzione comporta un'ulteriore perdita di pressione
totale, ma permette di limitare il numero di rampe e di ottenere una
configurazione geometricamente più compatta.

\paragraph{Dimensionamento geometrico.}
La geometria è stata costruita a partire dalla condizione teorica di
\textit{shock-on-lip}, facendo convergere gli urti verso il labbro della
presa. Rispetto alla configurazione ideale è stato tuttavia introdotto un
margine longitudinale $\Delta s$, arretrando il cowl rispetto al punto
teorico di convergenza. Tale margine è stato introdotto per rendere la
geometria meno sensibile alle variazioni del numero di Mach rispetto al
punto di progetto, che modificano l'inclinazione e quindi la posizione
degli urti. L'area di cattura $A_{\mathrm{catt}}$ è stata invece
determinata dalla portata d'aria richiesta dal ramjet mediante
l'equazione di continuità.

 Il procedimento di calcolo,
il confronto tra differenti numeri di urti e il dimensionamento
geometrico completo sono riportati in Appendice~\ref{app:presa_ramjet}.

La configurazione ottenuta presenta un recupero di pressione totale
$\pi_{\mathrm{tot}}=0.731$ e un'efficienza adiabatica
$\eta_d=0.879$. L'urto normale terminale porta il flusso a
$M_2=0.72$: tale valore, pur essendo subsonico, risulta ancora elevato
per l'ingresso nella camera di combustione e richiede pertanto
un'ulteriore diffusione subsonica a valle della presa.




% ============================================================
% PRESA D'ARIA TBCC - MACH 3.5
% ============================================================


% ============================================================
% 1. PARAMETRI PRINCIPALI
% ============================================================

% Dati ramjet.presaAria
\def\hCatt{1.1005}
\def\LSuper{2.2532}

\def\LRampaOneTwo{1.4907}
\def\LRampaTwoThree{0.7145}
\def\LRampaThreeFour{0.4435}

\def\wCatt{2.2011}

% Valori MATLAB
\def\deltaOneRaw{0.2057}
\def\deltaTwoRaw{0.2521}
\def\deltaThreeRaw{0.3093}

\def\thetaOneRaw{0.4544}
\def\thetaTwoRaw{0.5811}
\def\thetaThreeRaw{0.8046}

% Mach
\def\Mzero{3.5}

% Aree
\def\AtwoRam{0.5273}

% ATTENZIONE:
% sostituire col vero valore di motore.presaAria.A2
\def\AtwoTurbo{0.5273}

% Scala globale
\def\UnitScale{2.45}



% ============================================================
% 2. POSIZIONI MODIFICABILI DELLE SCRITTE
% ============================================================
%
% DX > 0  -> destra
% DX < 0  -> sinistra
%
% DY > 0  -> alto
% DY < 0  -> basso
% ============================================================

% theta 1
\def\ThetaOneDX{0.09}
\def\ThetaOneDY{0.18}

% delta 1
\def\DeltaOneDX{-0.03}
\def\DeltaOneDY{0.13}

% theta 2
\def\ThetaTwoDX{0.1}
\def\ThetaTwoDY{0.05}

% delta 2
\def\DeltaTwoDX{-0.02}
\def\DeltaTwoDY{0.14}

% theta 3
\def\ThetaThreeDX{0.1}
\def\ThetaThreeDY{0.1}

% delta 3
\def\DeltaThreeDX{-0.02}
\def\DeltaThreeDY{0.19}

% Shock point
\def\ShockTextDX{-0.5}
\def\ShockTextDY{0}

% Offset
\def\OffsetTextDX{-0.2}
\def\OffsetTextDY{-0.08}

% Ingresso compressore
\def\CompressorTextDX{-0.71}
\def\CompressorTextDY{-0.1}

% Mach
\def\MachTextDX{-0.10}
\def\MachTextDY{0.00}

% Post combustore
\def\PostTextDX{-0.2}
\def\PostTextDY{0.00}

% Etichette B C D E
\def\BTextDX{-0.05}
\def\BTextDY{0.05}

\def\CTextDX{0.00}
\def\CTextDY{0.06}

\def\DTextDX{0.00}
\def\DTextDY{0.06}

\def\ETextDX{0.00}
\def\ETextDY{0.06}



% ============================================================
% 3. CONVERSIONE ANGOLI
% ============================================================

\pgfmathsetmacro{\deltaOne}{deg(\deltaOneRaw)}
\pgfmathsetmacro{\deltaTwo}{deg(\deltaTwoRaw)}
\pgfmathsetmacro{\deltaThree}{deg(\deltaThreeRaw)}

\pgfmathsetmacro{\thetaOne}{deg(\thetaOneRaw)}
\pgfmathsetmacro{\thetaTwo}{deg(\thetaTwoRaw)}
\pgfmathsetmacro{\thetaThree}{deg(\thetaThreeRaw)}

\pgfmathsetmacro{\thetaFlow}{
    \deltaOne+\deltaTwo+\deltaThree
}



% ============================================================
% 4. AREE E ALTEZZE
% ============================================================

\pgfmathsetmacro{\hComp}{
    2*sqrt(\AtwoTurbo/pi)
}



% ============================================================
% 5. COORDINATE RAMPE
% ============================================================

\pgfmathsetmacro{\xB}{0}
\pgfmathsetmacro{\yB}{\hCatt}

\pgfmathsetmacro{\xC}{
    \xB+\LRampaOneTwo*cos(\deltaOne)
}

\pgfmathsetmacro{\yC}{
    \yB-\LRampaOneTwo*sin(\deltaOne)
}

\pgfmathsetmacro{\xD}{
    \xC+\LRampaTwoThree*cos(\deltaOne+\deltaTwo)
}

\pgfmathsetmacro{\yD}{
    \yC-\LRampaTwoThree*sin(\deltaOne+\deltaTwo)
}

\pgfmathsetmacro{\xE}{
    \xD+\LRampaThreeFour*cos(\thetaFlow)
}

\pgfmathsetmacro{\yE}{
    \yD-\LRampaThreeFour*sin(\thetaFlow)
}



% ============================================================
% 6. SHOCK POINT
% ============================================================

\pgfmathsetmacro{\xShock}{\LSuper}
\pgfmathsetmacro{\yShock}{0}



% ============================================================
% 7. OFFSET LABBRO
% ============================================================

\pgfmathsetmacro{\offsetLabbro}{
    0.05*\LSuper
}

\pgfmathsetmacro{\xLip}{
    \xShock+\offsetLabbro*cos(\thetaFlow)
}

\pgfmathsetmacro{\yLip}{
    \yShock-\offsetLabbro*sin(\thetaFlow)
}



% ============================================================
% 8. GEOMETRIA INTERNA
% ============================================================

\pgfmathsetmacro{\xHinge}{
    \xE+0.80*\LSuper
}

\pgfmathsetmacro{\xMax}{
    \xHinge+0.25*\LSuper
}

\pgfmathsetmacro{\yFloor}{
    -0.45*\hCatt
}

\pgfmathsetmacro{\yHinge}{
    \yFloor+sqrt(\AtwoRam/pi)
}

\pgfmathsetmacro{\yRoof}{
    \yHinge+\hComp
}



% ============================================================
% 9. PUNTI BEZIER - TETTO
% ============================================================

\pgfmathsetmacro{\RoofPoneX}{
    \xE+0.5*(\xHinge-\xE)
}

\pgfmathsetmacro{\RoofPoneY}{
    \yE
}

\pgfmathsetmacro{\RoofPtwoX}{
    \xHinge-0.2*(\xHinge-\xE)
}

\pgfmathsetmacro{\RoofPtwoY}{
    \yRoof
}



% ============================================================
% 10. PUNTI BEZIER - PARATIA MOBILE
% ============================================================

\pgfmathsetmacro{\dxFlap}{
    \xHinge-\xE
}

\pgfmathsetmacro{\FlapPoneX}{
    \xE+0.35*\dxFlap
}

\pgfmathsetmacro{\FlapPoneY}{
    \yE-0.25*\dxFlap*tan(\thetaFlow)
}

\pgfmathsetmacro{\FlapPtwoX}{
    \xHinge-0.35*\dxFlap
}

\pgfmathsetmacro{\FlapPtwoY}{
    \yHinge
}



% ============================================================
% 11. PUNTI BEZIER - PAVIMENTO RAMJET
% ============================================================

\pgfmathsetmacro{\dxFloor}{
    \xHinge-\xLip
}

\pgfmathsetmacro{\FloorPoneX}{
    \xLip+0.35*\dxFloor
}

\pgfmathsetmacro{\FloorPoneY}{
    \yLip-0.25*\dxFloor*tan(\thetaFlow)
}

\pgfmathsetmacro{\FloorPtwoX}{
    \xHinge-0.35*\dxFloor
}

\pgfmathsetmacro{\FloorPtwoY}{
    \yFloor
}



% ============================================================
% 12. ANGOLI ONDE D'URTO
% ============================================================

\pgfmathsetmacro{\aShockOne}{
    atan2(\yShock-\yB,\xShock-\xB)
}

\pgfmathsetmacro{\aShockTwo}{
    atan2(\yShock-\yC,\xShock-\xC)
}

\pgfmathsetmacro{\aShockThree}{
    atan2(\yShock-\yD,\xShock-\xD)
}



% ============================================================
% 13. DIREZIONE RAMPE
% ============================================================

\pgfmathsetmacro{\aRampOne}{
    -\deltaOne
}

\pgfmathsetmacro{\aRampTwo}{
    -(\deltaOne+\deltaTwo)
}

\pgfmathsetmacro{\aRampThree}{
    -(\deltaOne+\deltaTwo+\deltaThree)
}



% ============================================================
% 14. PARAMETRI ARCHI E FRECCE
% ============================================================

\pgfmathsetmacro{\rTheta}{
    0.16*\LSuper
}

\pgfmathsetmacro{\rDelta}{
    0.08*\LSuper
}

\pgfmathsetmacro{\Ldash}{
    0.20*\LSuper
}

\pgfmathsetmacro{\Larr}{
    0.08*\LSuper
}



% ============================================================
% 15. POSIZIONI FRECCE FLUSSO
% ============================================================

% Rampa 1
\pgfmathsetmacro{\mXOne}{
    (\xB+\xC)/2
}

\pgfmathsetmacro{\yShockOne}{
    \yB+
    (\yShock-\yB)/(\xShock-\xB)*
    (\mXOne-\xB)
}

\pgfmathsetmacro{\mYOne}{
    ((\yB+\yC)/2+\yShockOne)/2
}


% Rampa 2
\pgfmathsetmacro{\mXTwo}{
    (\xC+\xD)/2
}

\pgfmathsetmacro{\yShockTwo}{
    \yC+
    (\yShock-\yC)/(\xShock-\xC)*
    (\mXTwo-\xC)
}

\pgfmathsetmacro{\mYTwo}{
    ((\yC+\yD)/2+\yShockTwo)/2
}


% Rampa 3
\pgfmathsetmacro{\mXThree}{
    (\xD+\xE)/2
}

\pgfmathsetmacro{\yShockThree}{
    \yD+
    (\yShock-\yD)/(\xShock-\xD)*
    (\mXThree-\xD)
}

\pgfmathsetmacro{\mYThree}{
    ((\yD+\yE)/2+\yShockThree)/2
}


% Flusso interno
\pgfmathsetmacro{\mXFour}{
    (\xE+\xLip)/2+
    0.04*\LSuper*cos(\aRampThree)
}

\pgfmathsetmacro{\mYFour}{
    (\yE+\yLip)/2+
    0.04*\LSuper*sin(\aRampThree)
}



% ============================================================
% 16. FIGURA
% ============================================================

\begin{figure}[htbp]

\centering

\begin{tikzpicture}[
    scale = 0.8,
     transform shape,
    x=\UnitScale cm,
    y=\UnitScale cm,
    line cap=round,
    line join=round,
    font=\scriptsize,
    >={Stealth[length=2.6mm,width=2mm]}
]


% ============================================================
% FUSOLIERA SUPERIORE
% ============================================================

\draw[
    black,
    line width=0.8pt
]
({-0.15*\LSuper},\hCatt)
--
(\xMax,\hCatt);



% ============================================================
% RAMPE
% ============================================================

\draw[
    black,
    line width=1.3pt
]
(\xB,\yB)
--
(\xC,\yC)
--
(\xD,\yD)
--
(\xE,\yE);



% ============================================================
% ONDE D'URTO
% ============================================================

\draw[
    red,
    line width=1.3pt
]
(\xB,\yB)
--
(\xShock,\yShock);

\draw[
    red,
    line width=1.3pt
]
(\xC,\yC)
--
(\xShock,\yShock);

\draw[
    red,
    line width=1.3pt
]
(\xD,\yD)
--
(\xShock,\yShock);

\draw[
    red,
    line width=1.3pt
]
(\xE,\yE)
--
(\xShock,\yShock);



% ============================================================
% SHOCK POINT
% ============================================================

\fill[
    red
]
(\xShock,\yShock)
circle[radius=2pt];


\node[
    red,
    font=\bfseries\scriptsize
]
at
(
    {\xShock+\ShockTextDX},
    {\yShock+\ShockTextDY}
)
{
    Shock point
};



% ============================================================
% OFFSET LABBRO
% ============================================================

\draw[
    green!60!black,
    line width=3pt
]
(\xShock,\yShock)
--
(\xLip,\yLip);


\node[
    green!55!black,
    font=\bfseries\scriptsize
]
at
(
    {\xLip+\OffsetTextDX},
    {\yLip+\OffsetTextDY}
)
{
    Offset
};



% ============================================================
% TETTO TURBOGETTO
% ============================================================

\draw[
    black,
    line width=1.3pt
]
(\xE,\yE)
..
controls
    (\RoofPoneX,\RoofPoneY)
and
    (\RoofPtwoX,\RoofPtwoY)
..
(\xHinge,\yRoof)
--
(\xMax,\yRoof);



% ============================================================
% PARATIA MOBILE
% ============================================================

\draw[
    orange!85!black,
    line width=2.2pt
]
(\xE,\yE)
..
controls
    (\FlapPoneX,\FlapPoneY)
and
    (\FlapPtwoX,\FlapPtwoY)
..
(\xHinge,\yHinge);


\filldraw[
    fill=orange!85!black,
    draw=black,
    line width=0.5pt
]
(\xHinge,\yHinge)
circle[radius=2pt];



% ============================================================
% PAVIMENTO RAMJET
% ============================================================

\draw[
    black,
    line width=1.3pt
]
(\xLip,\yLip)
..
controls
    (\FloorPoneX,\FloorPoneY)
and
    (\FloorPtwoX,\FloorPtwoY)
..
(\xHinge,\yFloor)
--
(\xMax,\yFloor);



% ============================================================
% LINEA DIVISORIA
% ============================================================

\draw[
    black,
    line width=1.3pt
]
(\xHinge,\yHinge)
--
(\xMax,\yHinge);



% ============================================================
% INGRESSO COMPRESSORE
% ============================================================

\draw[
    magenta,
    line width=3pt
]
(\xMax,\yHinge)
--
(\xMax,\yRoof);


\node[
    magenta,
    font=\bfseries\scriptsize
]
at
(
    {\xMax+\CompressorTextDX},
    {(\yHinge+\yRoof)/2+\CompressorTextDY}
)
{
    Ingresso compressore
};



% ============================================================
% FRECCIA MACH
% ============================================================

\draw[
    ->,
    black,
    line width=1.5pt
]
({-0.1*\LSuper},{0.5*\hCatt})
--
({0.05*\LSuper},{0.5*\hCatt});


\node[
    font=\bfseries
]
at
(
    {-0.28*\LSuper+\MachTextDX},
    {0.5*\hCatt+\MachTextDY}
)
{
    $M_0=\Mzero$
};



% ============================================================
% POST COMBUSTORE
% ============================================================

\node[
    gray,
    font=\itshape\scriptsize
]
at
(
    {
      \xHinge
      +0.5*(\xMax-\xHinge)
      +\PostTextDX
    },
    {
      \yFloor
      +0.5*(\yHinge-\yFloor)
      +\PostTextDY
    }
)
{
    Post combustore $\rightarrow$
};



% ============================================================
% RAMPA 1
% ============================================================

\draw[
    dashed,
    cyan!70!blue,
    line width=0.8pt
]
(\xB,\yB)
--
++(\Ldash,0);


\draw[
    cyan!70!blue,
    line width=1pt
]
(\xB,\yB)
++(0:\rTheta)
arc[
    start angle=0,
    end angle=\aShockOne,
    radius=\rTheta
];


\draw[
    red,
    line width=1pt
]
(\xB,\yB)
++(0:\rDelta)
arc[
    start angle=0,
    end angle=\aRampOne,
    radius=\rDelta
];


\node[
    cyan!70!blue,
    font=\bfseries
]
at
(
    {
      \xB
      +1.12*\rTheta*cos(\aShockOne/2)
      +\ThetaOneDX
    },
    {
      \yB
      +1.12*\rTheta*sin(\aShockOne/2)
      +\ThetaOneDY
    }
)
{
    $\theta_1$
};


\node[
    red,
    font=\bfseries
]
at
(
    {
      \xB
      +1.30*\rDelta*cos(\aRampOne/2)
      +\DeltaOneDX
    },
    {
      \yB
      +1.30*\rDelta*sin(\aRampOne/2)
      +\DeltaOneDY
    }
)
{
    $\delta_1$
};


\draw[
    ->,
    black,
    line width=1.5pt
]
(
    {\mXOne-\Larr*cos(\aRampOne)/2+0.7},
    {\mYOne-\Larr*sin(\aRampOne)/2-0.35}
)
--
(
    {\mXOne+\Larr*cos(\aRampOne)/2+0.7},
    {\mYOne+\Larr*sin(\aRampOne)/2-0.35}
);



% ============================================================
% RAMPA 2
% ============================================================

\draw[
    dashed,
    cyan!70!blue,
    line width=0.8pt
]
(\xC,\yC)
--
++(\Ldash,0);


\draw[
    cyan!70!blue,
    line width=1pt
]
(\xC,\yC)
++(0:\rTheta)
arc[
    start angle=0,
    end angle=\aShockTwo,
    radius=\rTheta
];


\draw[
    red,
    line width=1pt
]
(\xC,\yC)
++(0:\rDelta)
arc[
    start angle=0,
    end angle=\aRampTwo,
    radius=\rDelta
];


\node[
    cyan!70!blue,
    font=\bfseries
]
at
(
    {
      \xC
      +1.12*\rTheta*cos(\aShockTwo/2)
      +\ThetaTwoDX
    },
    {
      \yC
      +1.12*\rTheta*sin(\aShockTwo/2)
      +\ThetaTwoDY
    }
)
{
    $\theta_2$
};


\node[
    red,
    font=\bfseries
]
at
(
    {
      \xC
      +1.30*\rDelta*cos(\aRampTwo/2)
      +\DeltaTwoDX
    },
    {
      \yC
      +1.30*\rDelta*sin(\aRampTwo/2)
      +\DeltaTwoDY
    }
)
{
    $\delta_2$
};


\draw[
    ->,
    black,
    line width=1.5pt
]
(
    {\mXTwo-\Larr*cos(\aRampTwo)/2+0.2},
    {\mYTwo-\Larr*sin(\aRampTwo)/2-0.2}
)
--
(
    {\mXTwo+\Larr*cos(\aRampTwo)/2+0.2},
    {\mYTwo+\Larr*sin(\aRampTwo)/2-0.2}
);



% ============================================================
% RAMPA 3
% ============================================================

\draw[
    dashed,
    cyan!70!blue,
    line width=0.8pt
]
(\xD,\yD)
--
++(\Ldash,0);


\draw[
    cyan!70!blue,
    line width=1pt
]
(\xD,\yD)
++(0:\rTheta)
arc[
    start angle=0,
    end angle=\aShockThree,
    radius=\rTheta
];


\draw[
    red,
    line width=1pt
]
(\xD,\yD)
++(0:\rDelta)
arc[
    start angle=0,
    end angle=\aRampThree,
    radius=\rDelta
];


\node[
    cyan!70!blue,
    font=\bfseries
]
at
(
    {
      \xD
      +1.12*\rTheta*cos(\aShockThree/2)
      +\ThetaThreeDX
    },
    {
      \yD
      +1.12*\rTheta*sin(\aShockThree/2)
      +\ThetaThreeDY
    }
)
{
    $\theta_3$
};


\node[
    red,
    font=\bfseries
]
at
(
    {
      \xD
      +1.30*\rDelta*cos(\aRampThree/2)
      +\DeltaThreeDX
    },
    {
      \yD
      +1.30*\rDelta*sin(\aRampThree/2)
      +\DeltaThreeDY
    }
)
{
    $\delta_3$
};


\draw[
    ->,
    black,
    line width=1.5pt
]
(
    {\mXThree-\Larr*cos(\aRampThree)/2},
    {\mYThree-\Larr*sin(\aRampThree)/2}
)
--
(
    {\mXThree+\Larr*cos(\aRampThree)/2},
    {\mYThree+\Larr*sin(\aRampThree)/2}
);



% ============================================================
% FRECCIA FLUSSO INTERNO
% ============================================================

\draw[
    ->,
    black,
    line width=1.5pt
]
(
    {\mXFour-\Larr*cos(\aRampThree)/2},
    {\mYFour-\Larr*sin(\aRampThree)/2}
)
--
(
    {\mXFour+\Larr*cos(\aRampThree)/2},
    {\mYFour+\Larr*sin(\aRampThree)/2}
);



% ============================================================
% ETICHETTE B C D E
% ============================================================

\node
at
(
    {\xB+\BTextDX},
    {\yB+\BTextDY}
)
{
    $B$
};


\node
at
(
    {\xC+\CTextDX},
    {\yC+\CTextDY}
)
{
    $C$
};


\node
at
(
    {\xD+\DTextDX},
    {\yD+\DTextDY}
)
{
    $D$
};


\node
at
(
    {\xE+\ETextDX},
    {\yE+\ETextDY}
)
{
    $E$
};

%%% lunghezaa


\end{tikzpicture}

\caption{
Presa d'aria TBCC a $M_0=\Mzero$.
}

\end{figure}





\begin{figure}[H]
    \centering

    % --- MINIPAGE DI SINISTRA: IMMAGINE ---
    \begin{minipage}[c]{0.55\textwidth}
        \centering
        % ============================================================
% PRESA D'ARIA TBCC - VERSIONE BIANCO E NERO
% ============================================================


% ============================================================
% 1. PARAMETRI PRINCIPALI
% ============================================================

% Dati reali MATLAB
\def\hCatt{1.1005}
\def\LSuper{2.2532}

\def\LRampaOneTwo{1.4907}
\def\LRampaTwoThree{0.7145}
\def\LRampaThreeFour{0.4435}

% Delta s
\def\DeltaS{0.2253}

% Angoli MATLAB originali
% Servono SOLO per costruire la geometria.
% Non vengono visualizzati.
\def\deltaOneRaw{0.2057}
\def\deltaTwoRaw{0.2521}
\def\deltaThreeRaw{0.3093}

% Mach
\def\Mzero{3.5}

% Aree
\def\AtwoRam{0.5273}

% Sostituire con motore.presaAria.A2 se hai il valore preciso
\def\AtwoTurbo{0.5273}

% Scala globale
\def\UnitScale{2.45}



% ============================================================
% 2. POSIZIONI MODIFICABILI DELLE SCRITTE
% ============================================================
%
% DX positivo = destra
% DX negativo = sinistra
%
% DY positivo = alto
% DY negativo = basso
% ============================================================

% Shock point
\def\ShockTextDX{-0.20}
\def\ShockTextDY{-0.08}

% Mach
\def\MachTextDX{-0.10}
\def\MachTextDY{0.00}

% Compressore
\def\CompressorTextDX{-0.15}
\def\CompressorTextDY{0.00}

% Post combustore
\def\PostTextDX{0.00}
\def\PostTextDY{0.00}

% B
\def\BTextDX{-0.04}
\def\BTextDY{0.06}

% C
\def\CTextDX{0.00}
\def\CTextDY{0.07}

% D
\def\DTextDX{0.00}
\def\DTextDY{0.07}

% E
\def\ETextDX{0.00}
\def\ETextDY{0.07}

% A = Shock point
\def\ATextDX{-0.08}
\def\ATextDY{0.04}

% A' = punto dopo Delta_s
\def\AprimeTextDX{0.05}
\def\AprimeTextDY{0.04}



% ============================================================
% 3. CONVERSIONE ANGOLI
% ============================================================

% I valori MATLAB vengono usati con sin/cos,
% quindi li convertiamo internamente per TikZ.

\pgfmathsetmacro{\deltaOne}{deg(\deltaOneRaw)}
\pgfmathsetmacro{\deltaTwo}{deg(\deltaTwoRaw)}
\pgfmathsetmacro{\deltaThree}{deg(\deltaThreeRaw)}

\pgfmathsetmacro{\thetaFlow}{
    \deltaOne+\deltaTwo+\deltaThree
}



% ============================================================
% 4. ALTEZZA COMPRESSORE
% ============================================================

\pgfmathsetmacro{\hComp}{
    2*sqrt(\AtwoTurbo/pi)
}



% ============================================================
% 5. COORDINATE RAMPE
% ============================================================

% B
\pgfmathsetmacro{\xB}{0}
\pgfmathsetmacro{\yB}{\hCatt}

% C
\pgfmathsetmacro{\xC}{
    \xB+\LRampaOneTwo*cos(\deltaOne)
}

\pgfmathsetmacro{\yC}{
    \yB-\LRampaOneTwo*sin(\deltaOne)
}

% D
\pgfmathsetmacro{\xD}{
    \xC+\LRampaTwoThree*cos(\deltaOne+\deltaTwo)
}

\pgfmathsetmacro{\yD}{
    \yC-\LRampaTwoThree*sin(\deltaOne+\deltaTwo)
}

% E
\pgfmathsetmacro{\xE}{
    \xD+\LRampaThreeFour*cos(\thetaFlow)
}

\pgfmathsetmacro{\yE}{
    \yD-\LRampaThreeFour*sin(\thetaFlow)
}



% ============================================================
% 6. SHOCK POINT = A
% ============================================================

\pgfmathsetmacro{\xShock}{\LSuper}
\pgfmathsetmacro{\yShock}{0}



% ============================================================
% 7. PUNTO A' DOPO DELTA S
% ============================================================

% Delta_s è interpretato come distanza orizzontale,
% come nella figura di riferimento.

\pgfmathsetmacro{\xLip}{
    \xShock+\DeltaS
}

% Il punto continua nella direzione della terza rampa
\pgfmathsetmacro{\yLip}{
    \yShock-\DeltaS*tan(\thetaFlow)
}



% ============================================================
% 8. GEOMETRIA INTERNA
% ============================================================

\pgfmathsetmacro{\xHinge}{
    \xE+0.80*\LSuper
}

\pgfmathsetmacro{\xMax}{
    \xHinge+0.25*\LSuper
}

\pgfmathsetmacro{\yFloor}{
    -0.45*\hCatt
}

\pgfmathsetmacro{\yHinge}{
    \yFloor+sqrt(\AtwoRam/pi)
}

\pgfmathsetmacro{\yRoof}{
    \yHinge+\hComp
}



% ============================================================
% 9. BEZIER - TETTO
% ============================================================

\pgfmathsetmacro{\RoofPoneX}{
    \xE+0.5*(\xHinge-\xE)
}

\pgfmathsetmacro{\RoofPoneY}{
    \yE
}

\pgfmathsetmacro{\RoofPtwoX}{
    \xHinge-0.2*(\xHinge-\xE)
}

\pgfmathsetmacro{\RoofPtwoY}{
    \yRoof
}



% ============================================================
% 10. BEZIER - PARATIA
% ============================================================

\pgfmathsetmacro{\dxFlap}{
    \xHinge-\xE
}

\pgfmathsetmacro{\FlapPoneX}{
    \xE+0.35*\dxFlap
}

\pgfmathsetmacro{\FlapPoneY}{
    \yE-0.25*\dxFlap*tan(\thetaFlow)
}

\pgfmathsetmacro{\FlapPtwoX}{
    \xHinge-0.35*\dxFlap
}

\pgfmathsetmacro{\FlapPtwoY}{
    \yHinge
}



% ============================================================
% 11. BEZIER - PAVIMENTO
% ============================================================

\pgfmathsetmacro{\dxFloor}{
    \xHinge-\xLip
}

\pgfmathsetmacro{\FloorPoneX}{
    \xLip+0.35*\dxFloor
}

\pgfmathsetmacro{\FloorPoneY}{
    \yLip-0.25*\dxFloor*tan(\thetaFlow)
}

\pgfmathsetmacro{\FloorPtwoX}{
    \xHinge-0.35*\dxFloor
}

\pgfmathsetmacro{\FloorPtwoY}{
    \yFloor
}



% ============================================================
% 12. FIGURA
% ============================================================


\centering

\begin{tikzpicture}[
    x=\UnitScale cm,
    y=\UnitScale cm,
    line cap=round,
    line join=round,
    font=\scriptsize,
    >={Stealth[length=2.5mm,width=1.8mm]}
]

\clip
  ({-0.42*\LSuper},{\yFloor-0.45})
  rectangle
  ({\xHinge-0.65*\LSuper},{\hCatt+0.15});

% ============================================================
% FUSOLIERA SUPERIORE / REFERENCE AXIS
% ============================================================

\draw[
    black,
    line width=0.9pt
]
({-0.15*\LSuper},\hCatt)
--
(\xMax,\hCatt);





% ============================================================
% RAMPE B-C-D-E
% ============================================================

\draw[
    black,
    line width=1.3pt
]
(\xB,\yB)
--
(\xC,\yC)
--
(\xD,\yD)
--
(\xE,\yE);



% ============================================================
% ONDE D'URTO - TUTTE NERE
% ============================================================

\draw[
    black,
    line width=1.1pt
]
(\xB,\yB)
--
(\xShock,\yShock);


\draw[
    black,
    line width=1.1pt
]
(\xC,\yC)
--
(\xShock,\yShock);


\draw[
    black,
    line width=1.1pt
]
(\xD,\yD)
--
(\xShock,\yShock);


\draw[
    black,
    line width=1.1pt
]
(\xE,\yE)
--
(\xShock,\yShock);



% ============================================================
% SHOCK POINT
% ============================================================




% ============================================================
% TRATTO A -> A'
% ============================================================

\draw[
    black,
    dashed,
    line width=1pt
]
(\xShock,\yShock)
--
(\xLip,\yLip);



% ============================================================
% ETICHETTE A E A'
% ============================================================

\node
at
(
    {\xShock+\ATextDX},
    {\yShock+\ATextDY-0.1}
)
{
    $A$
};


\node
at
(
    {\xLip+\AprimeTextDX},
    {\yLip+\AprimeTextDY}
)
{
    $A'$
};



% ============================================================
% TETTO
% ============================================================

\draw[
    black,
    line width=1.3pt
]
(\xE,\yE)
..
controls
    (\RoofPoneX,\RoofPoneY)
and
    (\RoofPtwoX,\RoofPtwoY)
..
(\xHinge,\yRoof)
--
(\xMax,\yRoof);



% ============================================================
% PARATIA MOBILE
% ============================================================

\draw[
    black,
    line width=1.6pt
]
(\xE,\yE)
..
controls
    (\FlapPoneX,\FlapPoneY)
and
    (\FlapPtwoX,\FlapPtwoY)
..
(\xHinge,\yHinge);


\fill[
    black
]
(\xHinge,\yHinge)
circle[radius=1.8pt];



% ============================================================
% PAVIMENTO
% ============================================================

\draw[
    black,
    line width=1.3pt
]
(\xLip,\yLip)
..
controls
    (\FloorPoneX,\FloorPoneY)
and
    (\FloorPtwoX,\FloorPtwoY)
..
(\xHinge,\yFloor)
--
(\xMax,\yFloor);



% ============================================================
% LINEA DIVISORIA
% ============================================================

\draw[
    black,
    line width=1.2pt
]
(\xHinge,\yHinge)
--
(\xMax,\yHinge);



% ============================================================
% FACCIA COMPRESSORE
% ============================================================

\draw[
    black,
    line width=2.3pt
]
(\xMax,\yHinge)
--
(\xMax,\yRoof);








% ============================================================
% ETICHETTE B C D E
% ============================================================

\node
at
(
    {\xB+\BTextDX},
    {\yB+\BTextDY}
)
{
    $B$
};


\node
at
(
    {\xC+\CTextDX},
    {\yC+\CTextDY}
)
{
    $C$
};


\node
at
(
    {\xD+\DTextDX},
    {\yD+\DTextDY}
)
{
    $D$
};


\node
at
(
    {\xE+\ETextDX},
    {\yE+\ETextDY}
)
{
    $E$
};



% ============================================================
% QUOTA VERTICALE h
% ============================================================

% Posizione della quota
\pgfmathsetmacro{\xHeightDim}{
    -0.38*\LSuper
}

% Linee di estensione
\draw[
    black,
    thin
]
(\xHeightDim,\hCatt)
--
(\xB,\hCatt);


\draw[
    black,
    thin
]
(\xHeightDim,\yShock)
--
(\xShock,\yShock);


% Freccia verticale
\draw[
    <->,
    black,
    line width=0.9pt
]
(\xHeightDim,\yShock)
--
(\xHeightDim,\hCatt);


% Etichetta h
\node[
    rotate=90,
    fill=white,
    inner sep=2pt
]
at
(
    \xHeightDim,
    {(\hCatt+\yShock)/2}
)
{
    $h$
};



% ============================================================
% QUOTA L_SUPER
% ============================================================

\pgfmathsetmacro{\yLsuperDim}{
    \yFloor-0.28
}

% Estensioni verticali
\draw[
    black,
    thin
]
(\xB,\yShock)
--
(\xB,\yLsuperDim);


\draw[
    black,
    thin
]
(\xShock,\yShock)
--
(\xShock,\yLsuperDim);


% Quota
\draw[
    <->,
    black,
    line width=0.9pt
]
(\xB,\yLsuperDim)
--
(\xShock,\yLsuperDim);


\node[
    fill=white,
    inner sep=2pt
]
at
(
    {(\xB+\xShock)/2},
    \yLsuperDim
)
{
    $L_{\mathrm{SUPER}}$
};



% ============================================================
% QUOTA DELTA_S
% ============================================================

\pgfmathsetmacro{\yDeltaDim}{
    \yLsuperDim+0.16
}

% Estensione shock point
\draw[
    black,
    thin
]
(\xShock,\yShock)
--
(\xShock,\yDeltaDim);


% Estensione A'
\draw[
    black,
    thin
]
(\xLip,\yLip)
--
(\xLip,\yDeltaDim);


% Quota Delta_s
\draw[
    <->,
    black,
    line width=0.9pt
]
(\xShock,\yDeltaDim)
--
(\xLip,\yDeltaDim);


\node[
    fill=white,
    inner sep=1pt
]
at
(
    {(\xShock+\xLip)/2},
    \yDeltaDim
)
{
    $\Delta_s$
};


\end{tikzpicture}


\caption{
Schema geometrico della presa d'aria TBCC.}

    \end{minipage}%
    \hfill
    % --- MINIPAGE DI DESTRA: TABELLA ---
    \begin{minipage}[c]{0.44\textwidth}
        \centering


               \captionof{table}{Parametri principali della presa d'aria supersonica del ramjet.}
               
        \scriptsize
        \renewcommand{\arraystretch}{1}
        \setlength{\tabcolsep}{3.5pt}

 
        \fontsize{11pt}{11pt}\selectfont
        \begin{tabular}{lcc}
            \toprule
            \textbf{Parametro} & \textbf{Valore} & \textbf{Unità} \\
            \midrule
            $\delta_1$ & 11.8 & $^\circ$ \\
            $\delta_2$ & 14.4 & $^\circ$ \\
            $\delta_3$ & 17.7 & $^\circ$ \\
            $\theta_1$ & 26.0 & $^\circ$ \\
            $\theta_2$ & 33.3 & $^\circ$ \\
            $\theta_3$ & 46.1 & $^\circ$ \\
            \midrule
            $\pi_{\mathrm{tot}}$ & 0.731 & -- \\
            $\eta_d$             & 0.879 & -- \\
            $M_2$                & 0.720 & -- \\
            $T_{02}$             & 765 & K \\
            $P_{02}$             & 142 & kPa \\
            \midrule
            $A_{\mathrm{catt}}$  & 2.42 & m$^2$ \\
            $A_2$                & 0.526 & m$^2$ \\
            $h_{\mathrm{catt}}$  & 1.10 & m \\
            $L_{\mathrm{super}}$ & 2.25 & m \\
            $\Delta s$           & 0.225 & m \\
            \bottomrule
        \end{tabular}

 
    \label{tab:miei_parametri}

        \label{tab:miei_parametri}
    \end{minipage}

\end{figure}


%-------------------------------------------------------------------------------------------------------------------------------------------------------























